\subsection{Warmup Prior Construction}
\label{sec:warmup_prior_construction}

A core architectural feature of banditGPT is its modular warmup prior pipeline: the system accepts \emph{any} collection of (prompt, model, reward) triples as input and produces LinUCB prior matrices ($\Amat, \bvec$) for each model in the portfolio.
Because warmup priors are inherently model-specific---they encode per-model reward surfaces over the prompt embedding space---the prior source must cover the models in the target portfolio.

\subsubsection{Same-Distribution Priors from LMSYS Arena}

All warmup priors in this paper are built from LMSYS Chatbot Arena prompts~\cite{zheng2023judging}, the same distribution used for online learning and holdout evaluation.
This \emph{same-distribution} design mirrors the most common production scenario: a practitioner collects a seed batch of representative workload prompts, evaluates them across all $K$ candidate models, and uses the resulting (prompt, model, reward) triples to initialise the router before it goes live.

Concretely, we reserve a 1{,}028-prompt \emph{prior-training} split from the LMSYS development data (strictly disjoint from the 1{,}543-prompt online-learning pool and the 1{,}483-prompt holdout; see Section~\ref{sec:three_way_split}).
The split is defined over the $K{=}10$ superset of models so that every portfolio ($K{=}2$, $K{=}3$, $K{=}10$) uses the same prompt partitions.
Each (prompt, model) pair contributes a rank-1 update to the LinUCB sufficient statistics: $\Amat_m \mathrel{+}= \xvec\xvec^\top$, $\bvec_m \mathrel{+}= r\xvec$.
A plasticity factor ($\rho{=}0.1$) scales the final matrices so that online observations can overwrite the prior within a practical number of steps.
For smaller portfolios ($K{=}2$, $K{=}3$), we filter the $K{=}10$ master artifact to retain only the relevant models, avoiding information leakage from models outside the test set.

\subsubsection{Three-Way Data Split for K\,$>$\,2}
\label{sec:three_way_split}

To support warmup prior construction without leaking information into the evaluation, we partition the 2{,}571 development prompts that have full $K{=}10$ model coverage into two disjoint subsets via stratified sampling (same category/complexity/difficulty axes as the canonical dev/holdout split, seed 42):

\begin{table}[htbp]
\centering
\caption{Three-way data split for all experiments.}
\label{tab:three_way_split}
\small
\resizebox{\columnwidth}{!}{%
\begin{tabular}{@{}llrl@{}}
\toprule
\textbf{Split} & \textbf{Source} & \textbf{Prompts} & \textbf{Purpose} \\
\midrule
Prior Training  & Dev ($K{=}10$ coverage) & 1{,}028 & Warmup priors for all 10 portfolio models \\
Online Learning & Dev ($K{=}10$ coverage) & 1{,}543 & Bandit online learning \\
Holdout         & Holdout ($K{=}10$ cov.) & 1{,}483 & Frozen evaluation (untouched) \\
\bottomrule
\end{tabular}
}
\end{table}

The prior-training set produces $1{,}028 \times 10 = 10{,}280$ (prompt, model, reward) observations---approximately 1{,}028 per model with a context dimension of 16, yielding ${\sim}64\times$ oversampling relative to the minimum needed for well-conditioned precision matrices.
This yields confident starting beliefs that online learning can still override within a few hundred steps thanks to the plasticity factor.

\paragraph{Data integrity.}
All three splits are verified prompt-disjoint by automated leakage checks at split creation time and again at experiment runtime.
The holdout set covers 1{,}483 prompts; experiments with smaller portfolios ($K{=}2$, $K{=}3$) read only the relevant model columns from the same holdout prompts, ensuring that comparisons across portfolio sizes are grounded in the same prompt population.

\subsubsection{Practical Guidance: Constructing Warmup Priors}
\label{sec:prior_guidance}

The warmup prior pipeline is intentionally generic: it requires only a set of (prompt, model, reward) triples in any format.
For practitioners deploying banditGPT with a new model portfolio, we recommend the following workflow:

\begin{enumerate}[nosep]
\item \textbf{Collect a seed batch.}
Evaluate 200--500 representative prompts from the target domain across all $K$ models in the portfolio.
Rewards can come from any source: human ratings, LLM-as-judge scores~\cite{zheng2023judging}, or automated quality metrics.
With $K$ models and $d{=}16$ context dimensions (15 PCA + bias), 200 prompts yield ${\sim}12\times$ oversampling per model, sufficient for confident priors.

\item \textbf{Generate priors.}
Run the prior generation script (\path{generate_multimodel_warmup_priors.py}), which encodes each prompt via the shared PCA projection and performs LinUCB updates for every (prompt, model, reward) triple.
The output is a single joblib artifact containing $\Amat_m$ and $\bvec_m$ for all $K$ models.

\item \textbf{Set plasticity.}
The plasticity factor $\rho$ (default 0.1) controls how quickly online observations can overwrite the prior.
Lower values ($\rho{=}0.01$) make the prior more persistent---appropriate when the seed batch closely matches deployment traffic.
Higher values ($\rho{=}1.0$) make the prior weaker---appropriate when the seed batch is from a different domain.
Section~\ref{sec:prior_degradation} shows that Corralling automatically hedges against prior quality, so practitioners need not precisely tune $\rho$.

\item \textbf{Deploy with Corralling.}
The Corralling meta-learner hedges between the warmup expert (initialized with the generated priors) and a tabula rasa expert (no priors).
If the warmup priors are well-calibrated, Corralling exploits them for 7--13\% lower regret than warmup-only (Section~\ref{sec:prior_degradation}).
If the priors are miscalibrated due to domain shift, Corralling detects the mismatch and shifts weight to the tabula rasa expert within ${\sim}50$ routing decisions.
\end{enumerate}

\noindent This workflow decouples the warmup prior from the router architecture: changing the model portfolio requires only regenerating the prior artifact, with no changes to the router code, hyperparameters, or deployment configuration.
All experiments in this paper use same-distribution priors from the LMSYS prior-training split, demonstrating the pipeline in the most common production scenario.
